% Options for packages loaded elsewhere
\PassOptionsToPackage{unicode}{hyperref}
\PassOptionsToPackage{hyphens}{url}
\documentclass[
]{article}
\usepackage{xcolor}
\usepackage{amsmath,amssymb}
\setcounter{secnumdepth}{-\maxdimen} % remove section numbering
\usepackage{iftex}
\ifPDFTeX
  \usepackage[T1]{fontenc}
  \usepackage[utf8]{inputenc}
  \usepackage{textcomp} % provide euro and other symbols
\else % if luatex or xetex
  \usepackage{unicode-math} % this also loads fontspec
  \defaultfontfeatures{Scale=MatchLowercase}
  \defaultfontfeatures[\rmfamily]{Ligatures=TeX,Scale=1}
\fi
\usepackage{lmodern}
\ifPDFTeX\else
  % xetex/luatex font selection
\fi
% Use upquote if available, for straight quotes in verbatim environments
\IfFileExists{upquote.sty}{\usepackage{upquote}}{}
\IfFileExists{microtype.sty}{% use microtype if available
  \usepackage[]{microtype}
  \UseMicrotypeSet[protrusion]{basicmath} % disable protrusion for tt fonts
}{}
\makeatletter
\@ifundefined{KOMAClassName}{% if non-KOMA class
  \IfFileExists{parskip.sty}{%
    \usepackage{parskip}
  }{% else
    \setlength{\parindent}{0pt}
    \setlength{\parskip}{6pt plus 2pt minus 1pt}}
}{% if KOMA class
  \KOMAoptions{parskip=half}}
\makeatother
\setlength{\emergencystretch}{3em} % prevent overfull lines
\providecommand{\tightlist}{%
  \setlength{\itemsep}{0pt}\setlength{\parskip}{0pt}}
\usepackage{bookmark}
\IfFileExists{xurl.sty}{\usepackage{xurl}}{} % add URL line breaks if available
\urlstyle{same}
\hypersetup{
  hidelinks,
  pdfcreator={LaTeX via pandoc}}

\author{}
\date{}

\begin{document}

\textbf{MedFinder -- Real-Time Medicine Availability \& Store Locator}

Project Proposal for UCS503P

Submitted to: Mr. Asif

\textbf{Thapar Institute of Engineering and Technology}

Project Team

Sourabh Chhabra(1024030435) \textbar{} Madhav Gupta(1024030436)
\textbar{} Arpit Yadav (1024030446)

August 16, 2026

\section{\texorpdfstring{\textbf{1. Higher-order goal}
\emph{\ldots secondary
framing}}{1. Higher-order goal \ldots secondary framing}}\label{higher-order-goal-secondary-framing}

This project aims to improve access to essential healthcare by reducing
the time patients spend searching for medicines during urgent or routine
needs. While the underlying issue of pharmacy inventory visibility spans
an entire city, the immediate engineering objective is to translate
real-time stock and location transparency into a measurable, deployable
software solution.

\section{\texorpdfstring{\textbf{2. Time-to-value} \emph{\ldots primary
heuristic}}{2. Time-to-value \ldots primary heuristic}}\label{time-to-value-primary-heuristic}

\begin{itemize}
\item
  We will deliver meaningful increments quickly by prototyping the core
  search-and-locate workflow and validating it with users within a short
  iteration window.
\item
  We will prioritize fast feedback over perfection by using weekly
  iterations (and targeting CI-backed automation early) to reduce
  release friction.
\item
  Success is defined by what can be delivered and measured in the first
  iteration --- a working medicine search and store locator --- then
  improved based on user feedback.
\end{itemize}

\section{\texorpdfstring{\textbf{3. Problem
Statement}}{3. Problem Statement}}\label{problem-statement}

Patients in urban and semi-urban areas frequently struggle to find
medicines quickly, especially during urgent medical needs. Common pain
points include:

\begin{itemize}
\item
  Fragmented availability: not every pharmacy stocks every medicine, and
  there is no way to know which store has a specific item without
  visiting or calling around.
\item
  Wasted time and effort: patients often roam from store to store, which
  delays treatment and worsens outcomes, particularly for the elderly or
  critically ill.
\item
  Unexpected stockouts: even common, everyday medicines go out of stock
  without warning, compounding the same problem.
\item
  No visibility before travel: there is no simple, centralized way to
  check medicine availability across nearby stores before stepping out
  of the house.
\end{itemize}

Why this matters now (contextual relevance): In cities where pharmacies
operate independently with no shared inventory system, even a
lightweight visibility layer can meaningfully cut down search time and
reduce delays in receiving treatment.

\section{\texorpdfstring{\textbf{4. Proposed Solution}
\emph{\ldots Software
Proposition}}{4. Proposed Solution \ldots Software Proposition}}\label{proposed-solution-software-proposition}

\subsection{\texorpdfstring{\textbf{4.1
Overview}}{4.1 Overview}}\label{overview}

Build a web-based ``Search-to-Pickup and Delivery'' platform with the
following components:

\begin{itemize}
\item
  Store \& inventory listing: recognized medical stores register and
  maintain a live inventory of the medicines they stock.
\item
  Medicine search: users search for a medicine and see which nearby
  stores have it in stock, along with price and distance.
\item
  Filtering \& sorting: results can be filtered/sorted by nearest
  distance and lowest price.
\item
  Reserve feature: users can place a short-hold reservation on a
  medicine at a chosen store so it is not sold out before they arrive.
\item
  Verified registration: only recognized, verified pharmacies can list
  inventory, keeping the platform trustworthy.
\end{itemize}

\subsection{\texorpdfstring{\textbf{4.2 Core
workflow}}{4.2 Core workflow}}\label{core-workflow}

\begin{itemize}
\item
  User searches for a medicine and instantly sees a ranked list of
  nearby stores with stock, price, and distance.
\item
  System checks store-submitted inventory data and flags low-stock or
  recently-updated listings.
\item
  User optionally reserves the medicine at a chosen store, receiving a
  reservation ID and hold window.
\item
  Store staff confirm the reservation, update stock upon pickup or
  expiry, and the record closes out.
\end{itemize}

\subsection{\texorpdfstring{\textbf{4.3 Operational constraints for an
educational
setting}}{4.3 Operational constraints for an educational setting}}\label{operational-constraints-for-an-educational-setting}

\begin{itemize}
\item
  Keep the solution usable on low-to-mid range devices via responsive
  web design.
\item
  Avoid dependence on advanced research algorithms (e.g., no ML-based
  demand forecasting in v1); focus on workflow, data correctness, and
  search performance.
\item
  Design for deployability using common web hosting and managed
  databases.
\end{itemize}

\section{\texorpdfstring{\textbf{5. Solution Approach}
\emph{\ldots Engineering
focus}}{5. Solution Approach \ldots Engineering focus}}\label{solution-approach-engineering-focus}

This is an engineering course project, so we focus on scalable
data-management and workflow aspects rather than novel algorithm
research.

\subsection{\texorpdfstring{\textbf{5.1 Search \& matching}
\emph{\ldots non-research}}{5.1 Search \& matching \ldots non-research}}\label{search-matching-non-research}

\begin{itemize}
\item
  Medicine search will use straightforward, well-understood techniques
  rather than novel research: fuzzy text matching on medicine
  name/brand/generic name to tolerate typos and partial queries.
\item
  Distance ranking will use standard geolocation math (haversine
  distance) against store coordinates, without claiming state-of-the-art
  routing.
\item
  Inventory freshness will be tracked via simple ``last updated''
  timestamps per store-medicine record, surfaced to users so stale
  listings are visible rather than hidden.
\end{itemize}

\subsection{\texorpdfstring{\textbf{5.2 Web architecture}
\emph{\ldots web-based
solution}}{5.2 Web architecture \ldots web-based solution}}\label{web-architecture-web-based-solution}

A typical 3-tier web architecture:

\begin{itemize}
\item
  Frontend: responsive UI for patients (search, filters, reservation)
  and store staff (inventory management dashboard).
\item
  Backend API: authentication, medicine search, inventory CRUD,
  reservation logic, store verification workflow.
\item
  Database: store users, stores, medicines, inventory records,
  reservations, and price history.
\end{itemize}

\subsection{\texorpdfstring{\textbf{5.3 CI/CD and fast delivery
enabler}}{5.3 CI/CD and fast delivery enabler}}\label{cicd-and-fast-delivery-enabler}

CI/CD will be used to:

\begin{itemize}
\item
  Automatically build and run tests on every change.
\item
  Produce repeatable deployment artifacts to staging.
\item
  Enable safe and frequent releases of incremental features (e.g.,
  adding the Reserve feature after core search is stable).
\end{itemize}

\section{\texorpdfstring{\textbf{6. Evaluation Criterion}
\emph{\ldots measurable and
attributable}}{6. Evaluation Criterion \ldots measurable and attributable}}\label{evaluation-criterion-measurable-and-attributable}

We define success using measurable metrics tied to the core workflow.

\subsection{\texorpdfstring{\textbf{6.1 Primary evaluation
metric}}{6.1 Primary evaluation metric}}\label{primary-evaluation-metric}

Time-to-find (TTF): time between a user submitting a medicine search
query and the app returning a usable list of in-stock nearby stores.

\begin{itemize}
\item
  Target: median TTF ≤ 3 seconds for search results during pilot window.
\item
  Attribution: measured via server-side request/response timestamps
  logged for each search query.
\end{itemize}

\subsection{\texorpdfstring{\textbf{6.2 Secondary evaluation
metrics}}{6.2 Secondary evaluation metrics}}\label{secondary-evaluation-metrics}

\begin{itemize}
\item
  Reservation fulfillment rate: percentage of reservations successfully
  honored (medicine still available) at pickup versus expired/cancelled.
\item
  Search success rate: percentage of searches that return at least one
  in-stock result, compared against a pre-launch baseline (phone-call
  survey of store availability).
\item
  User clarity: patient-reported clarity/usefulness of results using a
  simple 1--5 feedback question.
\item
  Reliability: availability of staging/production for search and
  reservation during business hours (e.g., ≥ 99\% uptime in pilot).
\end{itemize}

\subsection{\texorpdfstring{\textbf{6.3 Pilot validation plan}
\emph{\ldots high-level}}{6.3 Pilot validation plan \ldots high-level}}\label{pilot-validation-plan-high-level}

\begin{itemize}
\item
  Recruit 10--20 pilot users (patients) and 3--5 partner pharmacies (or
  simulated store accounts if real partners are unavailable).
\item
  Compare metrics before/after within the iteration window by logging
  baseline estimates via short interviews or manual store-calling
  exercises.
\end{itemize}

\section{\texorpdfstring{\textbf{7. Scalability} \emph{\ldots with
foresight}}{7. Scalability \ldots with foresight}}\label{scalability-with-foresight}

The proposition should scale in both theoretical and practical senses.

\subsection{\texorpdfstring{\textbf{Scalable
(theoretically)}}{Scalable (theoretically)}}\label{scalable-theoretically}

The system should support growth in number of stores, medicines, and
concurrent searches by:

\begin{itemize}
\item
  decoupling components (frontend/backend),
\item
  enabling pagination and indexing for common queries (e.g., medicine
  name, store location),
\item
  using stateless API design.
\end{itemize}

\subsection{\texorpdfstring{\textbf{Operationally deployable
(practically)}}{Operationally deployable (practically)}}\label{operationally-deployable-practically}

The system should run on common web hosting:

\begin{itemize}
\item
  managed database,
\item
  containerized or platform-hosted deployment,
\item
  CI/CD pipeline for staging/production.
\end{itemize}

\section{\texorpdfstring{\textbf{8. Engine availability
heuristic}}{8. Engine availability heuristic}}\label{engine-availability-heuristic}

A mature set of ``engines'' (tooling/services) is available to implement
this as a web-based project:

\begin{itemize}
\item
  Web framework for REST APIs and reactive frontend pages.
\item
  Managed relational or document database.
\item
  Authentication approach (token-based or session-based) with separate
  roles for patients and store staff.
\item
  Geolocation/maps service for distance calculation and map display.
\item
  Test framework for unit/integration tests.
\item
  CI runner and staging deployment target.
\end{itemize}

We will use these mature components to focus on workflow design,
correctness, and measurement rather than inventing new infrastructure.

\section{\texorpdfstring{\textbf{9. Project scope and deliverables}
\emph{\ldots iteration-friendly}}{9. Project scope and deliverables \ldots iteration-friendly}}\label{project-scope-and-deliverables-iteration-friendly}

\subsection{\texorpdfstring{\textbf{9.1 Initial deliverable}
\emph{\ldots first iteration}
}{9.1 Initial deliverable \ldots first iteration }}\label{initial-deliverable-first-iteration}

\begin{itemize}
\item
  Store registration and inventory submission form
\item
  Medicine search with distance and price filters
\item
  Patient results page: store list with stock status, price, distance
\item
  Basic logging and timestamps for TTF measurement
\item
  CI: build + backend unit tests + linting
\item
  CD to staging with a single command or pipeline job
\end{itemize}

\subsection{\texorpdfstring{\textbf{9.2 Subsequent
deliverables}}{9.2 Subsequent deliverables}}\label{subsequent-deliverables}

\begin{itemize}
\item
  Reserve feature with hold-timer and expiry handling
\item
  Notification layer (SMS/email optional in later iteration; can start
  with in-app notifications)
\item
  Store dashboard for inventory management and reservation metrics
\item
  Improved search/filtering and indexed queries for scale readiness
\end{itemize}

\section{\texorpdfstring{\textbf{10. Risks and
mitigations}}{10. Risks and mitigations}}\label{risks-and-mitigations}

\begin{itemize}
\item
  Data accuracy: store-submitted inventory may go stale; mitigate with
  ``last updated'' timestamps and periodic reminder prompts to stores.
\item
  Store adoption: provide an easy UI for inventory updates; keep the
  update workflow to a few taps; consider bulk-upload/CSV import.
\item
  Verification overhead: manual pharmacy verification could bottleneck
  onboarding; mitigate with a lightweight document-upload +
  admin-approval flow.
\item
  Evaluation measurement ambiguity: instrument timestamps and define
  clear search/reservation states before development begins.
\item
  Operational issues in pilot: use staging early; ensure CI/CD produces
  repeatable deployments.
\end{itemize}

\section{\texorpdfstring{\textbf{11.
Summary}}{11. Summary}}\label{summary}

This project proposes a deployable web platform that gives patients
real-time visibility into medicine availability across nearby
pharmacies. It meets the course criteria by emphasizing time-to-value
through rapid iterations, implementing CI/CD to support frequent safe
releases, and using measurable evaluation criteria (especially
time-to-find). It remains focused on engineering workflow, data
correctness, and scalability readiness rather than research-driven
novelty.

\end{document}
